\documentclass{article}

\usepackage{amssymb}
\usepackage{amsmath}

\title{Derivatives of the solution of the first order approximation of
a DSGE model}
\author{Michel Juillard}

\begin{document}
\maketitle

The first order approximation of a DSGE model is computed around the
deterministic steady state. Computing the derivatives of this
approximation with respect to the model parameters requires to compute
\begin{enumerate}
  \item the derivatives of the steady state with respect to
    the parameters.
  \item the derivatives of the coefficients of the linear
    approximation with respect to the parameters
  \item the derivatives of solution the Unilateral Quadratic
    Matrix Equation (UQME) used to obtain the solution of a linear
    rational expectation model
  \item the derivatives of the coeficients of response to shocks
\end{enumerate}

We consider the general model
\[
\mathbb{E}_t \left\{ f(y_{t+1}, y_t, y_{t-1}, \epsilon_t;
\theta)\right\} = 0
\]
where
\begin{align*}
  \mathbb{E}_t\left\{\epsilon_t\right\} &= 0\\
  \mathbb{E}_t\left\{\epsilon_t\epsilon_t'\right\} &= \Sigma_\epsilon\\
  \mathbb{E}_t\left\{\epsilon_t\epsilon_\tau'\right\} &= 0\;\;\;
  t\ne\tau
\end{align*}
$y_t$ is a vector of $n$ endogenous variables and $\epsilon_t$, a
vector of $m$ exogenous shocks.

\section{Derivatives of the steady state}

The steady state $\bar y$ satisfies
\[
f(\bar y, \bar y, \bar y, 0;\theta) = 0
\]
and
\[
\frac{\partial \bar y}{\partial \theta} = -\left(\frac{\partial
  f}{\partial y_{t+1}} +\frac{\partial f}{\partial y_t} + \frac{\partial
      f}{\partial y_{t-1}}\right)^{-1}\frac{\partial f}{\partial \theta}
\]

\section{Derivatives of the coefficents of the linear
  approximation of the model}  
The first order approximation of the model is written
\[
A \mathbb{E}_t\hat y_{t+1} + B \hat y_t + C \hat y_{t-1} + Du_t = 0
\]
where $\hat y_t=y_t - \bar y$ and
\begin{align*}
  A &= \left.\frac{\partial f}{\partial y_{t+1}} \right|_{\bar y}\\
  B &= \left.\frac{\partial f}{\partial y_t} \right|_{\bar y}\\
  C &= \left.\frac{\partial f}{\partial y_{t-1}} \right|_{\bar y}\\
  D &= \left.\frac{\partial f}{\partial u_t} \right|_{\bar y}
\end{align*}
Then,
\begin{align*}
  A &= \frac{\partial^2 f}{\partial y_{t+1}\partial \theta} + \left(\frac{\partial^2
  f}{\partial y_{t+1}\partial y_{t+1}} +\frac{\partial^2 f}{\partial y_t\partial y_{t+1}} + \frac{\partial^2
      f}{\partial y_{t-1}\partial y_{t+1}}\right)\left(e_n\otimes \frac{\partial \bar
    y}{\partial \theta}\right)\\ 
  B &= \frac{\partial^2 f}{\partial y_t\partial \theta} + \left(\frac{\partial^2
  f}{\partial y_{t+1}\partial y_t} +\frac{\partial^2 f}{\partial y_t\partial y_t} + \frac{\partial^2
      f}{\partial y_{t-1}\partial y_t}\right)\left(e_n\otimes \frac{\partial \bar
    y}{\partial \theta}\right)\\ 
  C &= \frac{\partial^2 f}{\partial y_{t-1}\partial \theta} + \left(\frac{\partial^2
  f}{\partial y_{t+1}\partial y_{t-1}} +\frac{\partial^2 f}{\partial y_t\partial y_{t-1}} + \frac{\partial^2
      f}{\partial y_{t-1}\partial y_{t-1}}\right)\left((e_n\otimes \frac{\partial \bar
    y}{\partial \theta}\right)\\ 
  D &= \frac{\partial^2 f}{\partial u_t\partial \theta} + \left(\frac{\partial^2
  f}{\partial y_{t+1}\partial u_t} +\frac{\partial^2 f}{\partial y_t\partial u_t} + \frac{\partial^2
      f}{\partial y_{t-1}\partial u_t}\right)\left(e_n\otimes \frac{\partial \bar
    y}{\partial \theta}\right)
\end{align*}
where $e_n$ is a 1-vector of size $n$.
\section{Derivatives of the solution of the UQME}

We assume that the solution of the first order approximation of the
model take the form
\[
\hat y_t = G\hat y_{t-1} + Hu_t
\]

Matrix $G$ must satisfies the UQME
\[
A G^2 + B G + C = 0
\]

We want to compute $\frac{\partial G}{\partial z}$ given $\frac{\partial
  A}{\partial z}$, $\frac{\partial B}{\partial z}$, $\frac{\partial
  C}{\partial z}$.

Derivative of an implicit function:
\[
  \frac{\partial A}{\partial z}G^2 + A\frac{\partial G}{\partial z}G
  + AG\frac{\partial G}{\partial z} + B\frac{\partial G}{\partial z}
  + \frac{\partial B}{\partial z}G + \frac{\partial C}{\partial z} = 0
\]
Can be reorganized as
\[
  A\frac{\partial G}{\partial z}G
  + (AG + B)\frac{\partial G}{\partial z} +\frac{\partial A}{\partial z}G^2 
  + \frac{\partial B}{\partial z}G + \frac{\partial C}{\partial z} = 0
\]
T
It has the form of a generalized Sylvester equation:
\[
  ax + bxc = d
\]
where, 
\begin{align*}
  a &= AG + B\\
  b &= A\\
  c &= G\\
  d &= \frac{\partial A}{\partial z}G^2 
      + \frac{\partial B}{\partial z}G + \frac{\partial C}{\partial z}\\
\end{align*}

This equation can be rewritten as
\[
  x + a^{-1}bxc = a^{-1}d,
\]  
and vectorized as
\[
  (I + c'\otimes(a^{-1}b)x = vec(a^{-1}d)
\]
The linear system can be solved for $x$.


\section{Derivatives of the coefficients of response to the shocks}
Matrix $H$ must satisfy
\[
A G H + BH + D = 0
\]
and
\[
H = -(AG + B)^{-1}D
\]

Then
\[
\frac{\partial H}{\partial \theta_i} = -(AG + B)\left(\frac{\partial
  A}{\partial \theta_i}GH + A\frac{\partial G}{\partial \theta_i}H +
  \frac{\partial B}{\partial \theta_i}H + D\right)
\]
\end{document}
